Departing from the previous research of the problem, a distinctive aspect of the model of this work is the generation of smooth advection fields as opposed to pixel-wise advection vector estimation.
The chosen approach is better suited especially for the high resolution meteorological applications in mind as small-scale differences in the data are mostly random given the available measurements, and features in the data mostly move to the same direction.

The data fed into the neural network consists of some smooth objects moving across two spatial dimensions over time.
Using several consecutive spatial snapshots of this data, the network predicts an advection field that represents the advection process across these time steps.
Even in the presence of noisy data, the objective is to produce smooth advectino fields.
The problem is depicted in figure~\ref{fig:advectionfigure}.
While the network is trained and evaluated mainly on satellite images of COT data, it is intended to be applicable to all similar processes.

Advection field estimation has also been studied in the context of optical flow problem.
Classical optical flow techniques are also often used in meteorology in nowcasting problems.
The neural network approach has also been studied but mostly focusing on different objectives.
They are often trained with ground truth flows and with datasets where motion only applies to small objects in images most notably with the KITTI~\cite{kitti}, and Sintel~\cite{sintel} datasets.
Contrary to this, the model used for this thesis has to be trained in an unsupervised fashion as the ground truth advection fields are not available, and the data for which the model is intended to be used with is very different from the usual ones used in the neural network optical flow context.

In order to tackle the needs of the much lower-dimensional data similar to figure~\ref{fig:advectionfigure}, a different kind of approach to the neural network structure is needed.
Basis vectors are utilised here akin to the method used in~\cite{deepide}, but a simple second degree polynomial is opted as opposed to a high-dimensional radial basis for the advection field.
The choice of outputting a low-dimensional polynomial basis from the neural network guarantees smooth advection fields. 

\subsubsection{Model Architecture}
The model architecture is similar to the well-known Residual Network (ResNet) that is deep yet computationally efficient.
The model compromises of two separate ResNet models, each trained to predict one component of the advection field: one for the vertical component, and the other for the horizontal component.

\subsubsection{ResNet Architecture}
Residual Networks are characterised by the use of skip connections, which allow the gradient to be directly backpropagated to earlier layers.
The residual term comes from these networks learning the residual mappings with reference to the layer inputs, alleviating the common vanishing gradient problem in deep neural networks.

A key component of the ResNet architecture is the residual block.
Residual blocks have the general form
\begin{equation}\label{eq:residualblock}
    F(x) + x,
\end{equation}
where $F(x)$ is the residual mapping to be learned representing layers of the neural network. $x$ is the original input to the block which represents the skip connection.
The important aspect here is the behaviour of the residual blocks in backpropagation.
Deep neural networks suffer from a problem called diminishing gradients when the chain rule results in a large number of multiplication of small numbers and thus producing small gradients.
However, differentiating the residual block gives
\begin{equation}\label{eq:residualblockderivative}
    \frac{d}{dx} F(x) + x = \frac{d}{dx} F(x) + 1.
\end{equation}
The $+1$ term ensures that even when propagating gradients through a deep network, they are not vanishing.
